\section{Cote 123: the tangent-rank prediction, Springer dimensions, and positivity}
\label{sec:cote123}

\subsection{The original letter and a corrected locator}
The original is Grothendieck's letter to Kostant dated 22 October 1969,
cote 123, physical PDF pages 22--23, manuscript and archive footers
21--22, in \url{https://grothendieck.umontpellier.fr/123.pdf}.
The official PDF has a cover. Thus the intake's locator ``123/023''
for the dimension question is a physical-PDF locator, not footer 23:
footer 23, physical page 24, is the nearly blank verso.
The entire two-page letter was inspected visually in this run.

On its second page Grothendieck asks how the tangent rank of the
adjoint quotient \(\chi:\mathfrak g\to\mathfrak h/W\), and of
the analogous group quotient, varies. For an orbit of dimension
\(\dim G-r-2\), he predicts rank \(r-1\) when all roots have
the same length, and strictly smaller rank otherwise. He also asks
for the dimension of the variety of Borel subgroups containing an
element, or whose Lie algebra contains a given Lie-algebra element.
Writing \(2N=\dim G-r\), an orbit of dimension \(2N-2i\)
is predicted to have such a variety of dimension \(i\).
The initial setting of the letter is a complex simple adjoint group.
These are mathematical questions in the letter itself, distinct
from the later intake's finite-field positivity interpretation.

\paragraph{Known answer to the dimension question.}
In the complex setting the asserted dimension is indeed
\[
 \dim\mathcal B_x=\frac{\dim Z_G(x)-r}{2}.
 \tag{123.1}
\]
Steinberg (1976), \textbf{Theorem 4.6, p.~217}, proves it for
unipotent group elements; \textbf{Theorem 6.5, p.~223}, proves
it for arbitrary group elements. His sufficient-characteristic
condition is automatic over \(\mathbf C\); it must not be
silently discarded when applying that paper in positive
characteristic. The Lie-algebra reduction is Springer's (1984)
\textbf{Lemma 1, p.~273}: split off the semisimple part and
pass to its connected reductive centralizer. In characteristic
zero a nilpotent element and its exponential have the same
centralizer and the same containing Borel subgroups. Applying
Steinberg's unipotent formula in that centralizer gives (123.1)
for the Lie-algebra question as well. This is an application
of known results, not a new proof of their geometry.

\subsection{BR-12301. The tangent rank in type A}
Over \(\mathbf C\), for \(X\in\mathfrak{gl}_n\), the tangent
rank of its adjoint quotient is the degree of its minimal
polynomial. For \(X\in\mathfrak{sl}_n\), it is that degree
minus one. In particular, a nilpotent matrix with largest
Jordan block of size \(b\) has ranks \(b\) and \(b-1\),
respectively. The same \(\mathfrak{sl}_n\) formula applies
to the Lie algebra of the adjoint group \(\mathrm{PGL}_n\).

\begin{proof}
The coefficients of the characteristic polynomial generate
the conjugation-invariant ring on \(\mathfrak{gl}_n\).
For clarity, one can check this assertion without choosing
generators from an abstract restriction theorem. Matrices
with distinct eigenvalues form a dense open subset; each is
conjugate to a diagonal matrix. The restriction of an
invariant polynomial is symmetric in the diagonal entries.
The elementary symmetric polynomials generate these symmetric
polynomials and lift to the characteristic coefficients.
Density proves that restriction is injective, which proves
the assertion. On trace-zero matrices the same argument uses
the hyperplane \(\sum_i\lambda_i=0\), and drops the first
characteristic coefficient. A polynomial on that hyperplane
invariant under the finite permutation group can be extended
invariantly to the whole diagonal space by averaging any
extension; hence its invariant ring is precisely the
symmetric-polynomial ring modulo the first coefficient.

In characteristic zero Newton identities replace the
characteristic coefficients, by a polynomial change of
coordinates with everywhere invertible triangular Jacobian,
with \(p_j(X)=\operatorname{Tr}(X^j)\), \(1\leq j\leq n\).
The differential is
\[
 dp_j|_X(Y)=j\operatorname{Tr}(X^{j-1}Y).
\]
The trace pairing on all matrices is nondegenerate, so their
span has dimension
\(\dim\langle 1,X,\ldots,X^{n-1}\rangle=\deg m_X\),
where \(m_X\) is the minimal polynomial. The last equality
is its defining minimality, together with division of
polynomials by \(m_X\).

The annihilator of \(\mathfrak{sl}_n\) for the trace pairing
is the scalar line: off-diagonal matrix units first force
a matrix in the annihilator to be diagonal, and differences
of diagonal units force all its entries to agree. The
quotient on \(\mathfrak{sl}_n\) uses \(p_2,\ldots,p_n\).
Their differentials therefore have rank equal to the
dimension of the image of \(\langle X,\ldots,X^{n-1}\rangle\)
modulo the scalar line. Together with that line these
powers span \(\mathbf C[X]\), so this rank is
\(\deg m_X-1\). For a nilpotent matrix \(m_X=t^b\).
Passing from \(\mathrm{SL}_n\) to its adjoint quotient
does not change the Lie algebra or its conjugation orbits.
This proves all assertions.
\end{proof}

\subsection{BR-12302. The strict nonsimply-laced rank prediction is false}
The strict inequality in the letter is refuted by the
complex simple adjoint group \(G=\mathrm{PSp}_4\), of type
\(C_2\). Put
\[
 J=\begin{pmatrix}0&I_2\\-I_2&0\end{pmatrix},\qquad
 N=\begin{pmatrix}0&I_2\\0&0\end{pmatrix}
       \in\mathfrak{sp}_4.
\]
Its orbit has dimension \(6=\dim G-r-2\), but
\(\operatorname{rank}(d\chi_N)=1=r-1\), despite the two
root lengths. The unipotent element represented by \(U=I+N\)
gives the same counterexample for the group quotient.
This refutes the stated strict inequality;
it does not assert a complete rank classification in every type.

\begin{proof}
Every element of \(\mathfrak{sp}_4\) has the form
\[
 Y=\begin{pmatrix}A&B\\C&-A^t\end{pmatrix},
 \qquad B=B^t,\quad C=C^t.
\]
Thus its dimension is \(4+3+3=10\) and its rank is two.
Direct multiplication gives
\[
 [Y,N]=\begin{pmatrix}-C&A+A^t\\0&C\end{pmatrix}.
\]
The Lie centralizer has \(C=0\), arbitrary symmetric
\(B\), and skew-symmetric \(A\), so dimension \(3+1=4\).
In characteristic zero an algebraic group stabilizer is
smooth; equivalently the conjugation orbit map has its
usual Lie-algebra tangent quotient. Therefore the orbit
has dimension \(10-4=6\). Passing through the finite
central isogeny to \(\mathrm{PSp}_4\) does not change it.

For completeness its adjoint quotient is explicitly
\(\chi(Y)=(c_2(Y),c_4(Y))\), where
\[
 \det(tI-Y)=t^4+c_2(Y)t^2+c_4(Y).
\]
Indeed, a regular semisimple symplectic matrix in the Lie
algebra has eigenvalues \(\lambda_1,\lambda_2,-\lambda_1,
-\lambda_2\). The symplectic form pairs opposite eigenspaces,
so an eigenbasis can be normalized to a symplectic basis.
Such matrices form a dense open subset. The normalizer
acts on the two \(\lambda_i\) by independent sign changes
and permutation. Its polynomial invariants are the
symmetric polynomials in \(\lambda_1^2,\lambda_2^2\),
generated by their sum and product. These lift to
\(-c_2\) and \(c_4\). Restriction and density prove
that these two coefficients generate the full invariant
ring, just as in the preceding proof.

Newton's identity gives \(c_2(Y)=-\tfrac12\operatorname{Tr}(Y^2)\).
Consequently \(dc_2|_N(Y)=-\operatorname{Tr}(NY)=-\operatorname{Tr}C\),
which is nonzero, for example when \(C=\operatorname{diag}(1,0)\).
Since \(N\) has rank two, all its three-by-three minors
vanish; the differential of \(c_4=\det\), whose entries
are those cofactors, is zero at \(N\). The quotient
differential thus has rank exactly one, as claimed.

Here is the group version, so that the letter's use of
unipotent orbits causes no ambiguity. Since \(N^2=0\) and
\(N^tJ=-JN\), direct expansion gives \(U^tJU=J\).
The centralizer of \(U=I+N\) is the centralizer of \(N\),
so its orbit also has dimension six. On a maximal torus
of \(\mathrm{Sp}_4\) the eigenvalues are
\(z_1,z_2,z_1^{-1},z_2^{-1}\). Put \(w_j=z_j+z_j^{-1}\).
Inversion-invariant Laurent polynomials in one variable
are polynomials in \(z+z^{-1}\), by the recurrence for
\(z^m+z^{-m}\). Independent inversions and permutation
therefore give invariant generators
\(t=w_1+w_2=\operatorname{Tr}g\) and
\(e=w_1w_2=c_2(g)-2\). Restriction to the dense regular
semisimple locus shows these generate the group invariant
ring, exactly as above. The central element \(-I\) sends
\((t,e)\) to \((-t,e)\); hence the adjoint group quotient
has coordinates \((t^2,e)\). At \(U\), \(t=4\), so
\((t,e)\mapsto(t^2,e)\) has invertible differential.

A tangent vector at \(U\) is \(UY\) with
\(Y\in\mathfrak{sp}_4\). Using
\(c_2(g)=((\operatorname{Tr}g)^2-\operatorname{Tr}(g^2))/2\),
\(U^2=I+2N\), and \(\operatorname{Tr}Y=0\), we obtain
\[
 dt_U(UY)=\operatorname{Tr}(NY),\qquad
 de_U(UY)=4\operatorname{Tr}(NY)-\operatorname{Tr}(U^2Y)
          =2\operatorname{Tr}(NY).
\]
The first functional is nonzero by the same choice of
\(C\), so the group quotient also has rank one, including
after descent to \(\mathrm{PSp}_4\).
\end{proof}

\subsection{BR-12303. The exact positive point-count consequence}
Let \(X_x\) be the Springer fibre of a nilpotent Lie-algebra
element \(x\) over \(\mathbf F_q\), with \(G\) connected reductive and residue
characteristic good for \(G\). Springer's (1984)
\textbf{Theorem 1, pp.~272--273}, says that \(H^j(X_x,\mathbf Q_\ell)\)
is pure of weight \(j\). His \textbf{Corollary 1 to Theorem 2,
p.~280}, gives odd vanishing, and \textbf{Lemma 4 on that page}
says that the eigenvalues on \(H^{2i}\) are \(q^i\) times
roots of unity. Here \(\ell\ne\operatorname{char}\mathbf F_q\).
These are the relevant finite-field inputs; purity alone
would not give the root-of-unity assertion.

There exists \(m\geq1\) such that, for all \(r\geq1\),
\[
 \#X_x(\mathbf F_{q^{mr}})
       =\sum_i b_{2i}q^{mri},\qquad
 Z(X_x/\mathbf F_{q^m},t)
       =\prod_i(1-q^{mi}t)^{-b_{2i}},
 \tag{123.2}
\]
where \(b_{2i}=\dim H^{2i}(X_x,\mathbf Q_\ell)\).
If all the root-of-unity factors are already one, one
can take \(m=1\). No Weyl-group trace is inserted in (123.2).

\begin{proof}
There are finitely many eigenvalues. Choose \(m\) divisible
by the orders of all their root-of-unity factors. The
eigenvalues of \(F^{mr}\) on \(H^{2i}\) are then all
\(q^{mri}\); this statement about traces needs no
semisimplicity assumption. The proper trace formula,
together with odd vanishing, proves the count. Substituting
it in \(Z(t)=\exp(\sum_{r\geq1}\#X_x(\mathbf F_{q^{mr}})t^r/r)\)
and using \(-\log(1-z)=\sum_{r\geq1}z^r/r\) proves the
displayed product as a formal power series. Every exponent
\(b_{2i}\) is nonnegative. This proves the stated consequence.
\end{proof}

\subsection{BR-12304. A Green function with a negative coefficient}
The intake SA20, \S4, D.018, goes from positive ordinary
counts to nonnegative coefficients of Green polynomials.
That implication is false. Use the Springer action and the
normalization of Springer (1984), \S5, p.~279:
\[
 Q_w(x)=\sum_i(-1)^i\operatorname{Tr}
       (Fw\mid H^i(X_x,\mathbf Q_\ell)).
 \tag{123.3}
\]
For \(G=\mathrm{SL}_2\), \(x=0\), odd \(q\), and the
nonidentity element \(s\in W=S_2\),
\[
 X_0=\mathbf P^1,\qquad Q_1(0)=1+q,
       \qquad Q_s(0)=1-q.
 \tag{123.4}
\]
The last polynomial has a negative coefficient. The
trace calculation holds for every odd \(q\). To invoke
the Green-function identification exactly as stated on
Springer's p.~279, take odd \(q\) sufficiently large;
the same polynomial is then a Green polynomial in that
source's terminology.

\begin{proof}
The regular semisimple orbit with eigenvalues \(1,-1\)
is the variety
\(U=(\mathbf P^1\times\mathbf P^1)\setminus\Delta\)
of ordered distinct eigenlines. The nonidentity Weyl element
interchanges the two lines. For \(x=0\), the simultaneous
resolution identifies its cohomology with that of \(X_0\);
this is precisely \textbf{Proposition 3 and the definition
of the Weyl action in \S5, p.~279}, of Springer (1984).
Thus it suffices to compute the swap action on \(H^*(U)\).

Projection \(U\to\mathbf P^1\) is a Zariski locally trivial
affine-line bundle: over a fixed line, complements to that
line form an affine line, and choosing a local complement
trivializes this torsor. Homotopy invariance gives cohomology
only in degrees zero and two, both of dimension one.
Let \(h_1,h_2\) be the two divisor classes on the product.
The divisor \(\Delta\) has class \(h_1+h_2\), as is seen
by intersecting it with each ruling. The localization
sequence for this smooth divisor yields
\[
 H^2(U,\mathbf Q_\ell)
   =\bigl(\mathbf Q_\ell h_1\oplus\mathbf Q_\ell h_2\bigr)
       /\mathbf Q_\ell(h_1+h_2).
\]
Here the adjacent \(H^1(\Delta)(-1)\) is zero and the
Gysin image is exactly the displayed divisor class.
The swap sends \(h_1\) to \(h_2=-h_1\), so it acts
by minus one on this one-dimensional group and by one
on \(H^0\). Both divisor classes are rational over
\(\mathbf F_q\), hence geometric Frobenius acts by
\(q\) on \(H^2\), and by one on \(H^0\).
Equation (123.3) now gives (123.4).
\end{proof}

The dimension formula and the ordinary positive count remain
valid. The failure occurs at the added Weyl trace, whose
character values need not be nonnegative. This is a concrete
boundary for any proposed passage from purity to positivity.
No RH or Dirichlet-GRH input is used in this section.

\subsection{Primary bibliography and inspection scope}
Grothendieck, A. (1969, October 22). \emph{Letter to B. Kostant}.
Montpellier archives, cote 123, physical pages 22--23,
footers 21--22. \url{https://grothendieck.umontpellier.fr/123.pdf}.
The letter also raises questions about opposed tori and rational
singularities; these are separate from the four statements proved
or cited above and are retained in the question register.

Springer, T. A. (1984). A purity result for fixed point varieties
in flag manifolds. \emph{Journal of the Faculty of Science,
University of Tokyo, Section IA, Mathematics, 31}(2), 271--282.
\url{https://doi.org/10.15083/00039529}.
Primary scan: \url{https://repository.dl.itc.u-tokyo.ac.jp/record/39538/files/jfs310202.pdf}.
Theorem 1, pp.~272--273; Lemma 1, p.~273; Proposition 3 and
the action convention, p.~279; Corollary 1 and Lemma 4, p.~280.
The stated results and the indicated action construction were read;
the full purity proof and its preceding geometric sources are
cited inputs, not independently reproved here.

Steinberg, R. (1976). On the desingularization of the unipotent
variety. \emph{Inventiones Mathematicae, 36}, 209--224.
\url{https://doi.org/10.1007/BF01390010}.
Theorem 4.6 and its characteristic convention, p.~217;
Remark 6.3 and Theorem 6.5, p.~223. Primary digitization:
\url{https://gdz.sub.uni-goettingen.de/download/pdf/PPN356556735_0036/LOG_0018.pdf}.
The complete short proofs of the two numbered statements were
read with their earlier classification and dimension results
retained as cited inputs.
